\subsection{A two-dimensional row module over a small local algebra}

\begin{lemma}[Proved: two-dimensional coefficients over a local algebra of dimension at most three]
\label{CP2:A-local-three-small-coefficient-complex}
Let $A$ be a rational algebra of dimension at most three with a unital
character $\chi:A\to\mathbb Q$ whose kernel $J$ is nilpotent.
Let $V$ be a left $A$-module of dimension at most two, and let $D$
be a complex of finite-dimensional right $A$-modules in degrees
zero and one.  If $D\otimes_A^{\mathbf L}V$ has bounded homology,
then its homology vanishes in every degree at least two.
\end{lemma}
\begin{proof}
The algebra is local with residue module $k={}_\chi\mathbb Q$.
The cases $V=0$ and $V$ projective are immediate.
If $JV=0$, then $V$ is a nonempty direct sum of copies of $k$,
so boundedness gives boundedness of $D\otimes_A^{\mathbf L}k$.
Otherwise $\dim V=2$, the radical $JV$ has dimension one,
and $V/JV$ has dimension one.  Nakayama's lemma supplies
a generator of $V$ and hence an exact sequence
\[
 0\longrightarrow U\longrightarrow A\longrightarrow V
 \longrightarrow0.
\]
Since $V$ is not projective, $\dim A$ cannot be two in this case:
the surjection would be an isomorphism.  Therefore $\dim A=3$
and $\dim U=1$.  Every one-dimensional module over this local
algebra is its residue module, so $U\cong k$.
Tensoring the sequence with $D$ gives
\[
 H_j(D\otimes_A^{\mathbf L}V)
   \cong H_{j-1}(D\otimes_A^{\mathbf L}k)
 \qquad(j\ge3),
\]
because $D\otimes_A^{\mathbf L}A\simeq D$ is two-term.
Thus boundedness again implies boundedness of the tensor with $k$.

In both remaining cases, take the minimal bounded-below free
resolution $P\to D$ constructed in the proof of
Lemma~\ref{CP2:A-local-two-term-derived-amplitude}.
All its differential entries lie in $J$, so $P\otimes_A k$
has zero differentials.  That lemma gives
$H_j(D\otimes_A^{\mathbf L}k)=0$ for $j\ge2$.
It follows that $P_j\otimes_A k=0$, hence $P_j=0$, for
$j\ge2$.  Consequently $D$ itself has a free resolution
in degrees zero and one.  Tensoring this resolution with
the original $V$ proves the assertion.  This last step is
needed at degree two; the displayed dimension shift alone
would only settle degrees at least three.
\end{proof}

\begin{theorem}[Proved: a local three-outside, three-row graph source]
\label{CP2:A-local-three-outside-row-graph}
Let $M$ have a proper rectangular minimal band with two columns
and at most three rows.  Put $A=\mathbb QM/\mathbb QK$, and
suppose $|M\setminus K|\le3$.
Suppose $A$ has a unital rational character with nilpotent kernel.
Let a finite nonempty left-$N$, right-$M$ biset have the natural
augmented graph-chain comparison and rationally acyclic opposite
corner of Theorem~\ref{CP2:A-two-column-two-outside-graph}.
If its collage has bounded rational homology, all rational
degree-two products vanish.
In particular the conclusion applies to $N=R_3$, with any state count.
\end{theorem}
\begin{proof}
Use the full two-column argument of
Theorem~\ref{CP2:A-two-column-two-outside-graph}.
The common-image-vertex case gives an actual bridge augmentation
section and kills products by the relative suspension argument.
That case uses no bound on the dimension of $A$.
For disjoint column-image graphs, the same injective projective
subcomplex $Ce\otimes eR$ and its two-term quotient $D$ give
\[
 H_n(Y;\mathbb Q)=H_{n-2}(D\otimes_A^{\mathbf L}V)
 \qquad(n\ge4),
\]
where $V$ is the full row-augmentation module and $Y$ is the
actual cofiber after collapsing the rationally acyclic opposite
corner.  None of the support-injectivity, augmentation-shift,
or projective-ideal change-of-rings arguments uses an outside
budget.  Here $\dim A\le3$ and $\dim V\le2$.
Boundedness of $Y$ gives boundedness of the displayed tensor.
The preceding lemma therefore kills its second homology,
so $H_4(Y;\mathbb Q)=0$.  The actual rational ring comparison
with the collage completes the proof.
\end{proof}
